\uuid{exQd}
\exo7id{3545}
\titre{exo7 3545}
\auteur{quercia}
\organisation{exo7}
\datecreate{2010-03-10}
\isIndication{false}
\isCorrection{false}
\chapitre{Réduction d'endomorphisme, polynôme annulateur}
\sousChapitre{Polynôme annulateur}
\cpge{mpsi-mp2i.applications-lineaires.b, pcsi.applications-lineaires.b, ptsi.applications-lineaires.b, mp-mpi.reduction-endomorphismes.b, pc.reduction.a, psi.reduction-endomorphismes.a, pt.reduction-endomorphismes.a, mp-mpi.reduction-endomorphismes.d, pc.reduction.c, psi.reduction-endomorphismes.c, pt.reduction-endomorphismes.c, mp-mpi.reduction-endomorphismes.i, pc.reduction.d, psi.reduction-endomorphismes.d, bcpst2.valeurs-propres-vecteurs-propres.a, bcpst2.valeurs-propres-vecteurs-propres.b}
\module{Algèbre}
\niveau{L2}
\difficulte{}

\contenu{
\texte{
Soit $E =  K_n[X]$ et $ u : E \to  E, P  \mapsto {X^nP(1/X).}$
}
\begin{enumerate}
    \item \question{Déterminer $u\circ u$. En déduire que $u$ est diagonalisable.}
    \item \question{Donner une base de vecteurs propres de $u$.}
\end{enumerate}
}
